% P15.29: the full Pólya source limit and its beta-mixture consequences.
\chapter{P\'olya Urns and Beta Prediction}
\label{sec:p15-polya-urns}

Repeated draws with reinforcement produce a random limiting proportion.
The finite-word probabilities determine its law and also give a concrete
model of updating and prediction. This chapter supplies the complete
urn presentation of Durrett's Section~4.3.2, including both special
cases and the general beta limit
\cite[pp.~226--227]{DurrettPTEAuthor2019}.
The beta-mixture construction below proves the general limit from its
already established integral normalization and the strong law.

\section{Construction and finite observations}

\begin{hbdefinition}{def:p15-polya-process}{Weights, draws and reinforcement}
Fix finite real \(g,r,c>0\). Initially the green and red weights are
\(g,r\). At each draw choose a color with probability proportional
to its current weight, retain the existing weights, and add \(c\)
to the chosen color. Integer parameters describe literal ball counts;
the same transition rule defines the stated real-weight model.

Let \(B_j\in\{0,1\}\) indicate a green draw, and put
\[
 S_0=0,\quad S_n=\sum_{j=1}^nB_j,\quad
 G_n=g+cS_n,\quad R_n=r+c(n-S_n),\quad
 X_n=\frac{G_n}{g+r+cn}.
\]
The time \(n\) refers to the state after \(n\) completed draws.
Write \(\mathcal F_n=\sigma(B_1,\ldots,B_n)\), with trivial
\(\mathcal F_0\), and
\[
                \alpha=g/c,\qquad\beta=r/c,\qquad t=\alpha+\beta.
\]
Thus \(X_n=(\alpha+S_n)/(t+n)\).
For \(a>0\) and integer \(k\geq0\), define the rising product
\((a)_0=1\), \((a)_k=\prod_{j=0}^{k-1}(a+j)\).
Empty products are one. A word \(b=(b_1,\ldots,b_n)\) has count
\(s=\sum_{j=1}^nb_j\).
\index{Pólya urn}\index{rising factorial}
\end{hbdefinition}

\begin{hbtheorem}{thm:p15-polya-construction}{proved=annex:p15.polya-construction}{Existence, the transition rule and bounded convergence}
Independent uniform variables \(U_j\) on \((0,1)\) give a construction
by \(B_{n+1}=\mathbf1_{\{U_{n+1}\leq X_n\}}\).
It has
\[
 P(B_{n+1}=1\mid\mathcal F_n)=X_n,\qquad
 X_{n+1}-X_n=\frac{B_{n+1}-X_n}{t+n+1}.
\]
Consequently \(X_n\) is a martingale in \([0,1]\), with
\(EX_n=\alpha/t\). It converges almost surely and in every
\(L^p\), \(0<p<\infty\), to a variable \(X_\infty\in[0,1]\).
Moreover \(S_n/n-X_n\to0\) for \(n\geq1\), so the draw frequency
has the same almost-sure limit.
\end{hbtheorem}

\begin{hbtheorem}{thm:p15-polya-word-law}{proved=annex:p15.polya-word-law}{The probability of every word and every count}
For every word \(b\) of length \(n\) and count \(s\),
\[
 P(B_1=b_1,\ldots,B_n=b_n)
       =\frac{(\alpha)_s(\beta)_{n-s}}{(t)_n}.
\]
In particular
\[
              P(S_n=s)=\binom ns
                    \frac{(\alpha)_s(\beta)_{n-s}}{(t)_n},
                    \quad 0\leq s\leq n.
\]
This is the beta-binomial law with parameters \(n,\alpha,\beta\).
The infinite draw sequence is exchangeable: its distribution is unchanged
by any permutation moving only finitely many coordinates.
Exchangeability does not assert independence.
\end{hbtheorem}

\section{The random limiting parameter}

\begin{hbdefinition}{def:p15-polya-mixture}{A beta parameter and conditional independent draws}
Let \(\Theta\) have the \(\operatorname{beta}(\alpha,\beta)\) law
defined in \cref{def:p15-beta-integral}, with density
\[
 \frac{\theta^{\alpha-1}(1-\theta)^{\beta-1}}{B(\alpha,\beta)}
                       \quad(0<\theta<1).
\]
Choose independent uniform \(V_j\) on \((0,1)\), mutually independent
and independent of \(\Theta\), and put
\(C_j=\mathbf1_{\{V_j\leq\Theta\}}\).
This is a joint construction on a product probability space.
For each fixed parameter value \(\theta\), its coordinate section
is an independent Bernoulli sequence of success probability \(\theta\).
No independence of the unconditioned \(C_j\) is asserted.
\end{hbdefinition}

\begin{hbtheorem}{thm:p15-polya-mixture-representation}{proved=annex:p15.polya-mixture-representation}{Equality of the full draw laws}
For a word \(b\) with count \(s\), and every Borel \(D\subseteq[0,1]\),
\[
 P(\Theta\in D,C_1=b_1,\ldots,C_n=b_n)
 =\frac{1}{B(\alpha,\beta)}
       \int_D\theta^{\alpha+s-1}(1-\theta)^{\beta+n-s-1}\,d\theta,
\]
where the integrand is zero at the endpoints.
The sequence \(C\) has the same law as \(B\) on
\(\{0,1\}^{\mathbb N}\), including all events in its product sigma-field.
Also \(n^{-1}\sum_{j=1}^nC_j\to\Theta\) almost surely.
The proof uses a measurable convergence event and product integration,
rather than assuming that parameter-dependent null sets coincide.
\end{hbtheorem}

\begin{hbtheorem}{thm:p15-polya-beta-limit}{proved=annex:p15.polya-beta-limit}{The general source limit and its joint meaning}
Write \(Z\sim\eta\) to mean that the random variable \(Z\) has law \(\eta\).
For every \(g,r,c>0\),
\[
       X_\infty\sim\operatorname{beta}(g/c,r/c),\qquad
       S_n/n\longrightarrow X_\infty\quad\text{almost surely}.
\]
For each fixed finite word, the joint law of that word and \(X_\infty\)
is the joint law given in \cref{thm:p15-polya-mixture-representation}
with \(C,\Theta\) replaced by \(B,X_\infty\).
In particular, conditionally on \(X_\infty\), the draws have the
Bernoulli product law in the sense of every finite joint conditional
probability. This is an exact probability representation of the urn
law; it does not change the physical reinforcement rule.
\end{hbtheorem}

\begin{hbtheorem}{thm:p15-polya-posterior-prediction}{proved=annex:p15.polya-posterior-prediction}{Updating the limiting fraction and predicting several draws}
After a word of length \(n\) and count \(s\), put
\(a=\alpha+s\), \(b=\beta+n-s\), \(T=a+b=t+n\).
The conditional law of \(X_\infty\) is
\(\operatorname{beta}(a,b)\), with
\[
 E(X_\infty\mid\mathcal F_n)=X_n,\qquad
 \operatorname{Var}(X_\infty\mid\mathcal F_n)
          =\frac{X_n(1-X_n)}{t+n+1}.
\]
For integer \(m\geq1\), let \(K_{n,m}=\sum_{j=1}^mB_{n+j}\).
Its conditional probabilities are
\[
 P(K_{n,m}=k\mid\mathcal F_n)
       =\binom mk\frac{(a)_k(b)_{m-k}}{(T)_m},
                    \quad 0\leq k\leq m.
\]
In particular
\[
 E(K_{n,m}\mid\mathcal F_n)=\frac{ma}{T},\qquad
 \operatorname{Var}(K_{n,m}\mid\mathcal F_n)
       =\frac{mab(T+m)}{T^2(T+1)}.
\]
For \(m>1\) this variance exceeds the binomial variance
\(mX_n(1-X_n)\). The several-draw predictive law retains uncertainty
about the common limiting fraction.
\end{hbtheorem}

\begin{hbtheorem}{thm:p15-polya-moments}{proved=annex:p15.polya-moments}{Dependence, finite-count variation and prediction error}
For nonnegative integers \(k,l\),
\[
 E[X_\infty^k(1-X_\infty)^l]
            =\frac{(\alpha)_k(\beta)_l}{(t)_{k+l}}.
\]
In particular \(EX_\infty=\alpha/t\) and
\(\operatorname{Var}(X_\infty)=\alpha\beta/[t^2(t+1)]\).
For distinct \(i,j\),
\[
 \operatorname{Cov}(B_i,B_j)=\frac{\alpha\beta}{t^2(t+1)},
 \qquad \operatorname{Corr}(B_i,B_j)=\frac1{t+1}.
\]
At every \(n\geq0\),
\[
 \begin{aligned}
 ES_n&=\frac{n\alpha}{t},&
 \operatorname{Var}(S_n)&=\frac{n\alpha\beta(t+n)}{t^2(t+1)},\\
 \operatorname{Var}(X_n)&=\frac{n\alpha\beta}{t^2(t+1)(t+n)},&
 E(X_\infty-X_n)^2&=\frac{\alpha\beta}{t(t+1)(t+n)}.
 \end{aligned}
\]
The positive covariance explicitly distinguishes this exchangeable
sequence from independent draws with fixed success probability \(\alpha/t\).
\end{hbtheorem}

\section{Source calculations and parameter boundaries}

\begin{hbtheorem}{thm:p15-polya-source-special-cases}{proved=annex:p15.polya-source-special-cases}{Uniform and linear beta densities}
If \(g=r=c=1\), \(S_n\) is uniform on \(\{0,\ldots,n\}\),
\(G_n=S_n+1\), and \(X_\infty\) is uniform on \((0,1)\).

If \(g=2,r=c=1\), then
\[
 P(G_n=s+2)=\frac{2(s+1)}{(n+1)(n+2)},\qquad
 P(S_n\leq k)=\frac{(k+1)(k+2)}{(n+1)(n+2)}
\]
for \(0\leq s,k\leq n\). The limiting fraction has density \(2x\)
on \(0<x<1\) and CDF \(x^2\) on \([0,1]\).
For integers \(0\leq s_n\leq n\) with \(s_n/n\to x\),
\[
             nP(G_n=s_n+2)\longrightarrow2x,
             \qquad P(G_n=s_n+2)\longrightarrow0.
\]
The single-cell probability and the limiting density are distinct.
The annex retains the source's direct finite-mass route to both limits.
\end{hbtheorem}

\begin{hbtheorem}{thm:p15-polya-parameter-boundaries}{proved=annex:p15.polya-parameter-boundaries}{Reinforcement strength and degenerate parameters}
Multiplying all three positive parameters by the same positive constant
leaves the draw law unchanged.
For fixed \(g,r>0\), put \(p=g/(g+r)\).
As \(c\downarrow0\), the beta limit law converges weakly to \(\delta_p\).
As \(c\to\infty\), it converges weakly to
\((1-p)\delta_0+p\delta_1\).

For the separately defined zero-reinforcement process \(c=0\),
the fraction is constantly \(p\), the draws are independent Bernoulli
variables of probability \(p\), and \(S_n/n\to p\) almost surely.
The positive-\(c\) beta parameters are not evaluated at \(c=0\).
For \(c\geq0\), if \(g=0<r\), the process draws only red almost
surely and its fraction limit has law \(\delta_0\).
If \(r=0<g\), it draws only green almost surely and the limit
has law \(\delta_1\).
With both initial weights zero, the first proportional draw is undefined.
\end{hbtheorem}

\begin{hbtheorem}{thm:p15-polya-worked-prediction}{proved=annex:p15.polya-worked-prediction}{An ordered observation, future counts and model evidence}
Start with \(g=2,r=c=1\), and observe the ordered word green, green, red.
Its probability is \(1/10\); the probability of two green draws in
any order is \(3/10\).
The updated limiting fraction has law \(\operatorname{beta}(4,2)\),
mean \(2/3\), and variance \(2/63\).
The count of green draws in the next two draws has probabilities
\[
                    1/7,\quad 8/21,\quad 10/21
\]
at \(0,1,2\), respectively, and variance \(32/63\).
In particular the probability of two green draws is \(10/21\),
rather than the independent plug-in value \(4/9\).

For model comparison, let \(H_U\) be this urn model and \(H_F\)
the model of independent fair draws. Define the Bayes factor for
\(H_U\) against \(H_F\) at the ordered data to be the ratio of their
data probabilities. It equals \(4/5\).
Positive finite prior odds \(O\) for \(H_U\) become posterior odds
\(4O/5\); with equal prior model probabilities,
\(P(H_U\mid\text{green, green, red})=4/9\).
The conditional prediction inside \(H_U\) and the probability assigned
to \(H_U\) are different inferential quantities.
\end{hbtheorem}

\begin{proofboundary}{def:p15-polya-source-scope}{The completed urn presentation and subsequent exchangeability}
The entire Section~4.3.2 source presentation is supplied: transition
martingale, every finite word/count probability, both displayed special
cases and the beta limit for every positive parameter triple.
The beta-mixture proof and the updating, dependence and parameter-limit
calculations are supplied additions. The opposite-color reinforcement
model has its distinct conclusion in \cref{thm:p15-friedman-urn-limit}.
The general conditional-product theorem is
\cref{thm:p15-definetti-factorization}; the probability-valued
representation on standard Borel spaces is
\cref{thm:p15-definetti-kernel}. The urn's specific beta mixing law
here gives a computed instance of those general results.
\end{proofboundary}
